\section{Insert Interval}
\label{sec:insert-interval}

\problemstatement{We are given a list of intervals $\textit{intervals}$ that
is already sorted by start time and contains no overlaps (it represents,
say, a calendar of existing non-conflicting events). We are also given one
new interval $\textit{newInterval}$. We must insert
$\textit{newInterval}$ into the list, merging it with any existing
intervals it overlaps, so that the result is still a list of intervals
sorted by start time with no overlaps.}

\begin{patternbox}
\emph{``The existing intervals are already sorted and non-overlapping; insert
one new interval and merge where necessary''} is a strong hint that we do
not need to re-sort or re-scan everything from scratch -- the existing order
is a resource we should exploit with a single linear pass, split into three
natural phases (intervals entirely before the new one, intervals overlapping
the new one, intervals entirely after). Whenever a problem promises that the
input is already sorted, a greedy single pass that respects that order is
almost always available, replacing what would otherwise look like a
sort-and-merge operation.
\end{patternbox}

\subsection*{Understanding the problem carefully}

Because the existing intervals are sorted by start time and mutually
non-overlapping, as we scan through them from left to right, each interval
falls into exactly one of three categories relative to
$\textit{newInterval} = [\textit{ns}, \textit{ne}]$:

\begin{enumerate}
  \item \textbf{Entirely before}: the interval's end is strictly less than
        $\textit{ns}$. It cannot overlap $\textit{newInterval}$ at all, and
        because the list is sorted, every interval we have not yet visited
        could still potentially overlap, but every interval already
        confirmed to end before $\textit{ns}$ can be output immediately,
        unchanged.
  \item \textbf{Overlapping}: the interval's start is at most
        $\textit{ne}$ and its end is at least $\textit{ns}$. Any such
        interval must be merged into a single growing interval together
        with $\textit{newInterval}$ and every other overlapping interval.
  \item \textbf{Entirely after}: the interval's start is strictly greater
        than $\textit{ne}$. Once we reach such an interval, we know
        $\textit{newInterval}$ (possibly already merged with some earlier
        intervals) is finished growing, because the remaining intervals,
        being sorted, can only start even later.
\end{enumerate}

\begin{ideabox}[Greedy Choice]
Scan the existing sorted intervals once, left to right. First, copy every
interval that ends strictly before $\textit{newInterval}$ starts directly
into the answer, untouched. Next, while the current interval overlaps the
(possibly already-growing) new interval, absorb it by extending the new
interval's start to the minimum of the two starts and its end to the maximum
of the two ends. Once no more intervals overlap, insert this final merged
interval into the answer. Finally, copy every remaining interval directly
into the answer, untouched.
\end{ideabox}

\subsection*{Why a single pass suffices}

This works cleanly because of two facts guaranteed by the problem: the
existing intervals are sorted by start time, and they do not overlap each
other. Sortedness means that once we pass an interval ending before
$\textit{ns}$, no later interval can also need to be classified as
``entirely before'' incorrectly -- the boundary between ``before'' and
``possibly overlapping'' is crossed exactly once as we scan forward.
Non-overlap among the existing intervals means that, when merging the
overlapping group, we never need to worry about two existing intervals
needing to be merged with each other independently of
$\textit{newInterval}$ -- they only ever come into contact with each other
\emph{through} their shared overlap with the growing merged interval. This
is what allows the entire operation to be a single $O(n)$ pass with no
need to sort or to revisit any interval more than once.

\subsection*{Algorithm}

Write $[s_i, e_i]$ for $\textit{intervals}[i]$, and $[\textit{ns},
\textit{ne}]$ for the (possibly growing) new interval.

\begin{enumerate}
  \item Initialize an empty result list and an index $i \gets 0$.
  \item \textbf{Phase 1.} While $i < n$ and $e_i < \textit{ns}$: append
        $[s_i, e_i]$ to the result; $i \gets i+1$.
  \item \textbf{Phase 2.} While $i < n$ and $s_i \le \textit{ne}$: merge by
        setting $\textit{ns} \gets \min(\textit{ns}, s_i)$ and
        $\textit{ne} \gets \max(\textit{ne}, e_i)$; $i \gets i+1$. Once the
        loop ends, append the final merged $[\textit{ns}, \textit{ne}]$ to
        the result.
  \item \textbf{Phase 3.} While $i < n$: append $[s_i, e_i]$ to the result;
        $i \gets i+1$.
  \item Return the result.
\end{enumerate}

\begin{algorithm}[H]
\caption{Insert Interval}
\begin{algorithmic}[1]
\Procedure{InsertInterval}{\textit{intervals}, \textit{newInterval}}
  \State $\textit{result} \gets$ empty list, $i \gets 0$, $n \gets \text{length}(\textit{intervals})$
  \While{$i < n$ \textbf{and} $\textit{intervals}[i].\textit{end} < \textit{newInterval}.\textit{start}$}
    \State append $\textit{intervals}[i]$ to \textit{result}; $i \gets i+1$
  \EndWhile
  \While{$i < n$ \textbf{and} $\textit{intervals}[i].\textit{start} \le \textit{newInterval}.\textit{end}$}
    \State $\textit{newInterval}.\textit{start} \gets \min(\textit{newInterval}.\textit{start}, \textit{intervals}[i].\textit{start})$
    \State $\textit{newInterval}.\textit{end} \gets \max(\textit{newInterval}.\textit{end}, \textit{intervals}[i].\textit{end})$
    \State $i \gets i + 1$
  \EndWhile
  \State append \textit{newInterval} to \textit{result}
  \While{$i < n$}
    \State append $\textit{intervals}[i]$ to \textit{result}; $i \gets i+1$
  \EndWhile
  \State \Return \textit{result}
\EndProcedure
\end{algorithmic}
\end{algorithm}

\subsection*{Dry run}

\dryrun{Take $\textit{intervals} = [[1,3],[6,9]]$ and
$\textit{newInterval} = [2,5]$.}

\begin{itemize}
  \item \textbf{Phase 1}: check $[1,3]$. Is $3 < 2$? No. Phase 1 ends with
        nothing copied, $i=0$.
  \item \textbf{Phase 2}: check $[1,3]$. Is $1 \le 5$? Yes, overlaps. Merge:
        $\textit{newInterval} = [\min(2,1), \max(5,3)] = [1,5]$. $i=1$.
        Check $[6,9]$. Is $6 \le 5$? No. Phase 2 ends. Append $[1,5]$ to
        result.
  \item \textbf{Phase 3}: append $[6,9]$ unchanged.
\end{itemize}

The result is $[[1,5],[6,9]]$.

Now take $\textit{intervals} = [[1,2],[3,5],[6,7],[8,10],[12,16]]$ and
$\textit{newInterval} = [4,8]$.

\begin{itemize}
  \item \textbf{Phase 1}: $[1,2]$ has end $2 < 4$, copy it. $[3,5]$ has end
        $5 \not< 4$, stop. Result so far: $[[1,2]]$, $i$ points at $[3,5]$.
  \item \textbf{Phase 2}: $[3,5]$: start $3 \le 8$, overlaps. Merge:
        $\textit{newInterval}=[\min(4,3),\max(8,5)]=[3,8]$. Next, $[6,7]$:
        start $6\le8$, overlaps. Merge: $[\min(3,6),\max(8,7)]=[3,8]$
        (unchanged). Next, $[8,10]$: start $8\le8$, overlaps. Merge:
        $[\min(3,8),\max(8,10)]=[3,10]$. Next, $[12,16]$: start $12\le10$?
        No, stop. Append $[3,10]$ to result.
  \item \textbf{Phase 3}: append $[12,16]$ unchanged.
\end{itemize}

The result is $[[1,2],[3,10],[12,16]]$.

\subsection*{C++ implementation}

\begin{lstlisting}
#include <bits/stdc++.h>
using namespace std;

vector<vector<int>> insertInterval(vector<vector<int>>& intervals,
                                    vector<int>& newInterval) {
    vector<vector<int>> result;
    int n = intervals.size();
    int i = 0;

    // Phase 1: intervals ending before newInterval starts
    while (i < n && intervals[i][1] < newInterval[0]) {
        result.push_back(intervals[i]);
        i++;
    }

    // Phase 2: merge all overlapping intervals
    while (i < n && intervals[i][0] <= newInterval[1]) {
        newInterval[0] = min(newInterval[0], intervals[i][0]);
        newInterval[1] = max(newInterval[1], intervals[i][1]);
        i++;
    }
    result.push_back(newInterval);

    // Phase 3: remaining intervals
    while (i < n) {
        result.push_back(intervals[i]);
        i++;
    }

    return result;
}

int main() {
    vector<vector<int>> intervals = {{1,2},{3,5},{6,7},{8,10},{12,16}};
    vector<int> newInterval = {4, 8};

    auto result = insertInterval(intervals, newInterval);
    for (auto& iv : result) {
        cout << "[" << iv[0] << "," << iv[1] << "] ";
    }
    cout << endl;
    return 0;
}
\end{lstlisting}

\begin{complexitybox}
\textbf{Time Complexity.} Each of the $n$ existing intervals is visited
exactly once across the three phases combined (no interval is revisited),
and each visit does constant work, giving
\[
  O(n).
\]
\textbf{Space Complexity.} The result list holds at most $n+1$ intervals, so
the space complexity is $O(n)$ for the output (no other significant
auxiliary storage is used).
\end{complexitybox}

\begin{takeawaybox}
When the input is explicitly guaranteed to already be sorted, look for a
single linear pass with a small, fixed number of phases (here:
before-overlap-after) instead of reapplying a general sort-and-merge
algorithm. Exploiting a sortedness guarantee given for free in the problem
statement is itself a greedy-style optimization: it converts an
$O(n \log n)$ general merge into an $O(n)$ specialized one.
\end{takeawaybox}
